\chapter{Samples, Concentration, and Limit Theorems}
\chaptermark{Samples, Concentration, and Limit Theorems}
\label{chap:limit-theorems}

\begingroup
\clubpenalty=10000
\widowpenalty=10000

\section{An average changes from dataset to dataset}
\label{sec:samples}

Record one hundred waiting times from the inbox of
Section~\ref{sec:exponential} and report their average. Repeat the
experiment with another hundred waits. The two averages will usually
differ. How far might either average be from the population mean,
and how does collecting more observations change that uncertainty?

Keep the same model: independent waits
\(T_i\sim\operatorname{Exp}(0.2)\), measured in minutes. Their mean
and standard deviation are both five minutes, as calculated in
Chapter~\ref{chap:distributions}. Here the new object is the
\emph{dataset} and the average computed from it, not another distribution
for an individual wait.

\subsection{A population, a sample, and a sampling distribution}

The \emph{population distribution} describes the observations we want
to learn about. A \emph{sample} is the collection
\(X_1,\ldots,X_n\) that we observe. We call these observations
\emph{independent and identically distributed}, abbreviated i.i.d.,
when they are mutually independent and all follow the same distribution.
Both requirements matter: measurements from different operating
conditions need not share a distribution, and repeated measurements
of one object need not be independent.

Before collection, \(X_1,\ldots,X_n\) are random variables; after
collection, \(x_1,\ldots,x_n\) are the recorded numbers. The
\emph{sample mean} is
\begin{equation}
  \bar X_n=\frac1n\sum_{i=1}^n X_i.
  \label{eq:sample-mean}
\end{equation}
Its distribution over repetitions of the entire sampling experiment
is a \emph{sampling distribution}. It is not the distribution of
values within one dataset. To visualize the distinction, imagine
collecting many independent datasets of size one hundred, calculating
one average for each, and plotting those averages.

For i.i.d. observations with finite mean \(\mu\) and variance
\(\sigma^2\), the sum and scaling rules from
Sections~\ref{sec:expectation} and~\ref{sec:variance} immediately give
\begin{equation}
  \Expect[\bar X_n]=\mu,\qquad
  \Var(\bar X_n)=\frac{\sigma^2}{n}.
  \label{eq:sample-mean-moments}
\end{equation}
This is the general form of the click-proportion calculation in
Section~\ref{sec:bernoulli-binomial}. We now use it to study accuracy.

\subsection{Standard deviation versus standard error}

The standard deviation of a statistic across repeated samples is
called its \emph{standard error}. For an i.i.d. sample mean,
\begin{equation}
  \operatorname{SE}(\bar X_n)=\frac{\sigma}{\sqrt n}.
  \label{eq:standard-error}
\end{equation}
For the inbox, a single wait has standard deviation five minutes,
while the average of one hundred waits has standard error
\(5/\sqrt{100}=0.5\) minutes. Individual waits remain just as variable
when the dataset grows. It is their average that becomes more precise.

In this particular model, the Gamma sum result in
Section~\ref{sec:gamma} supplies an exact sampling distribution:
\begin{equation}
  \sum_{i=1}^n T_i\sim\operatorname{Gamma}(n,0.2),\qquad
  \bar T_n\sim\operatorname{Gamma}(n,0.2n).
  \label{eq:exponential-sample-mean}
\end{equation}
The second parameter is a rate. Dividing a positive variable by \(n\)
multiplies its Gamma rate by \(n\). Figure~\ref{fig:ch4-sampling}
shows how the resulting density narrows while retaining mean five.

\begin{figure}[H]
  \centering
  \begin{tikzpicture}
\begin{axis}[
  width=0.95\textwidth,height=6.6cm,
  axis lines=left,xmin=0,xmax=16,ymin=0,ymax=0.46,
  xlabel={Average waiting time (minutes)},ylabel={Density},
  xtick={0,2,4,6,8,10,12,14,16},ytick={0,0.1,0.2,0.3,0.4},
  tick label style={font=\small},label style={font=\small},
  grid=major,grid style={black!8},
  legend style={draw=none,font=\small,at={(0.5,1.04)},anchor=south,
    legend columns=3,column sep=1.3em},
]
  \addplot[CourseBlue,very thick] table[x=x,y=n1] {figures/data/ch4-sampling-density.dat};
  \addlegendentry{\(n=1\)}
  \addplot[orange!85!black,very thick,dashed] table[x=x,y=n4] {figures/data/ch4-sampling-density.dat};
  \addlegendentry{\(n=4\)}
  \addplot[teal!80!black,very thick,densely dotted] table[x=x,y=n25] {figures/data/ch4-sampling-density.dat};
  \addlegendentry{\(n=25\)}
  \addplot[CourseGray,thin,dashed,forget plot] coordinates {(5,0) (5,0.435)};
  \node[font=\small,anchor=west,text=CourseDark,fill=white,inner sep=3pt]
    at (axis cs:7.2,0.35) {Same mean: 5 minutes};
\end{axis}
\end{tikzpicture}

  \caption{Exact sampling densities of the mean of \(n\) independent
  exponential waits, each with mean five minutes. For \(n=1\), the
  curve is the distribution of a single observation. The larger-sample
  curves describe averages of entire datasets, not less variable
  individual waits. Their standard errors are \(5\), \(2.5\), and
  \(1\) minutes.}
  \label{fig:ch4-sampling}
\end{figure}

In a real dataset, the population variance is usually unknown. A common
measure of the observed spread is the \emph{sample variance},
\[
  S_n^2=\frac1{n-1}\sum_{i=1}^n(X_i-\bar X_n)^2,
  \qquad n\ge2.
\]
The denominator \(n-1\) corrects for centering the data at their own
sample mean; the statistical-estimation chapter will derive that
correction. Writing \(S_n=\sqrt{S_n^2}\), the corresponding estimate
of the mean's standard error is \(S_n/\sqrt n\). It is an estimate,
not the exact value of \(\sigma/\sqrt n\), and its interpretation
depends on the sampling assumptions.

\section{Bounding the probability of a large error}
\label{sec:concentration}

Suppose we want the average waiting time to be within one minute of
the population mean. Knowing its standard error gives a scale for
the error, but does not by itself give a probability. We first seek a
bound that does not require the full sampling distribution.

\subsection{Markov: a nonnegative quantity cannot often be very large}

Let \(W\ge0\) have finite expectation. On the event \(W\ge a\),
where \(a>0\), its contribution is at least \(a\). Pointwise,
\[
  W\ge a\,\mathbf1_{\{W\ge a\}}.
\]
Taking expectations, using the indicator rule from
Section~\ref{sec:expectation}, gives
\(\Expect[W]\ge a\Prob(W\ge a)\). Rearranging yields
\begin{theorem}[Markov's inequality]
For \(W\ge0\), \(\Expect[W]<\infty\), and \(a>0\),
\begin{equation}
  \Prob(W\ge a)\le\frac{\Expect[W]}{a}.
  \label{eq:markov}
\end{equation}
\end{theorem}

For a nonnegative wait with mean five minutes, Markov gives
\(\Prob(T\ge20)\le1/4\). The exponential model gives the much smaller
exact value \(e^{-4}\approx0.0183\). Markov is valid for \emph{every}
nonnegative distribution with that mean, including one that puts
probability \(1/4\) at twenty and \(3/4\) at zero. For that latter
distribution the bound is attained. Without more information about
the distribution, we cannot improve this bound uniformly.

Nonnegativity is essential. A signed quantity could take very large
positive and negative values whose contributions cancel in its mean.
To control deviations on both sides, apply Markov to their squares.

\subsection{Chebyshev: turn a variance into an error bound}

For \(X\) with mean \(\mu\) and finite variance \(\sigma^2\), take
\(W=(X-\mu)^2\) and \(a=\varepsilon^2\) in Markov's inequality.
The event \(W\ge\varepsilon^2\) is exactly
\(|X-\mu|\ge\varepsilon\). Thus
\begin{theorem}[Chebyshev's inequality]
For \(\varepsilon>0\),
\begin{equation}
  \Prob(|X-\mu|\ge\varepsilon)
  \le\frac{\sigma^2}{\varepsilon^2}.
  \label{eq:chebyshev}
\end{equation}
\end{theorem}

Applying this to \(\bar X_n\), using
\eqref{eq:sample-mean-moments}, gives
\begin{equation}
  \Prob(|\bar X_n-\mu|\ge\varepsilon)
  \le\min\!\left(1,\frac{\sigma^2}{n\varepsilon^2}\right).
  \label{eq:chebyshev-mean}
\end{equation}
The minimum with one simply discards a bound larger than any probability.
For one hundred inbox waits and tolerance one minute,
\[
  \Prob(|\bar T_{100}-5|\ge1)\le\frac{25}{100}=0.25.
\]
This is an upper bound, not a prediction that a quarter of datasets
will miss the target.

Conversely, a sufficient sample size for failure probability at most
\(\delta\in(0,1)\) is
\[
  n\ge\frac{\sigma^2}{\delta\varepsilon^2}.
\]
For the inbox, \(n=500\) guarantees an error below one minute except
on an event of probability at most \(0.05\). The guarantee uses the
true variance, or a known upper bound on it. Substituting a random
sample variance without further justification does not preserve the
same finite-sample guarantee.

\section{The law of large numbers}
\label{sec:lln}

Chebyshev answers a question at a specified sample size. It also gives
a limiting conclusion: for a fixed error tolerance, the probability
of exceeding it tends to zero as the number of observations grows.

\begin{theorem}[Weak law of large numbers: finite-variance version]
Let \(X_1,X_2,\ldots\) be i.i.d. with mean \(\mu\) and finite
variance. Then, for every \(\varepsilon>0\),
\begin{equation}
  \lim_{n\to\infty}\Prob(|\bar X_n-\mu|>\varepsilon)=0.
  \label{eq:weak-lln}
\end{equation}
\end{theorem}
\begin{proof}
By \eqref{eq:chebyshev-mean}, the probability is bounded above by
\(\sigma^2/(n\varepsilon^2)\), which tends to zero.
\end{proof}

This statement is called \emph{convergence in probability}, written
\(\bar X_n\xrightarrow{\Prob}\mu\). Choose a fixed tolerance first,
such as one minute or one tenth of a minute. Among repeatedly collected
datasets, the fraction whose average lies outside that tolerance becomes
arbitrarily small as the dataset size increases.

\begin{figure}[H]
  \centering
  \begin{tikzpicture}
\begin{axis}[
  width=0.95\textwidth,height=6.7cm,
  axis lines=left,xmode=log,log basis x=10,xmin=1,xmax=1000,ymin=0,ymax=12,
  xlabel={Number of observations \(n\) (log scale)},ylabel={Running average (minutes)},
  xtick={1,10,100,1000},xticklabels={1,10,100,1000},ytick={0,2,4,6,8,10,12},
  tick label style={font=\small},label style={font=\small},
  grid=major,grid style={black!8},
  legend style={draw=none,font=\small,at={(0.5,1.04)},anchor=south,
    legend columns=3,column sep=0.9em},
]
  \addplot[CourseBlue,thick] table[x=n,y=path1] {figures/data/ch4-running-means.dat};
  \addlegendentry{Dataset 1}
  \addplot[orange!85!black,thick,dashed] table[x=n,y=path2] {figures/data/ch4-running-means.dat};
  \addlegendentry{Dataset 2}
  \addplot[teal!80!black,thick,dashdotted] table[x=n,y=path3] {figures/data/ch4-running-means.dat};
  \addlegendentry{Dataset 3}
  \addplot[CourseGray,thin,densely dashed,forget plot] coordinates {(1,5) (1000,5)};
  \node[font=\small,text=CourseDark,anchor=north east,fill=white,inner sep=3pt]
    at (axis cs:900,4.8) {True mean: 5};
\end{axis}
\end{tikzpicture}

  \caption{Three independent simulated sequences of exponential waits.
  Each curve uses the first \(n\) observations of its own sequence.
  The horizontal axis is logarithmic. Averages can move away from the
  mean after a new observation; convergence does not mean monotone
  improvement. These are illustrative paths, not a probability bound.}
  \label{fig:ch4-running}
\end{figure}

The theorem does not require the error to decrease at every step of
one running experiment. A new large wait can move an otherwise accurate
average farther from five. Nor does it claim that individual waits
approach five: their distribution is unchanged. The running averages
in Figure~\ref{fig:ch4-running} fluctuate even as their overall spread
shrinks.

Finite variance is sufficient for this short proof, but is not necessary
for the i.i.d. law of large numbers. The more general result requires
only \(\Expect[|X_1|]<\infty\). That distinction will matter for the
heavy-tailed examples in Section~\ref{sec:averaging-failures}.

The distribution being sampled also matters. If only short waits are
recorded, their average can settle very precisely near the mean of that
selected population, not the mean of all waits. Increasing sample size
reduces sampling fluctuations; it does not correct an unrepresentative
sampling process.

\section{The central limit theorem}
\label{sec:clt}

The law of large numbers tells us that the error becomes small. It does
not tell us its shape. For the inbox we have an exact Gamma answer;
for a complicated observation or a model's loss, that distribution may
be unavailable. The central limit theorem supplies a broadly applicable
approximation under stated assumptions.

\subsection{Keep the shrinking fluctuations visible}

The error \(\bar X_n-\mu\) has standard deviation
\(\sigma/\sqrt n\). Divide by this scale:
\[
  Z_n=\frac{\bar X_n-\mu}{\sigma/\sqrt n}
     =\frac{\sum_{i=1}^n(X_i-\mu)}{\sigma\sqrt n}.
\]
This expresses the sampling error in standard-error units. The error
itself shrinks, while \(Z_n\) has mean zero and variance one for every
\(n\). Equal mean and variance do not determine a shape; the new
claim concerns the distribution of \(Z_n\).

\begin{theorem}[Classical i.i.d. central limit theorem]
Let \(X_1,X_2,\ldots\) be i.i.d. with mean \(\mu\) and
\(0<\sigma^2=\Var(X_1)<\infty\). Then
\begin{equation}
  \frac{\sqrt n(\bar X_n-\mu)}{\sigma}
  \xrightarrow{d}\mathcal N(0,1).
  \label{eq:clt}
\end{equation}
Explicitly, for every real \(z\),
\[
  \Prob(Z_n\le z)\longrightarrow\Phi(z).
\]
\end{theorem}

The symbol \(\xrightarrow{d}\) means \emph{convergence in distribution}:
here, convergence of the CDFs to the Gaussian CDF. This is different
from saying that the rescaled errors converge to zero. They retain
nonzero spread. It also does not say that the original \(X_i\) become
Gaussian or that the unscaled sample mean approaches a Gaussian of
fixed positive variance.

Returning to the original units gives the approximation
\[
  \bar X_n\ \approx\ \mathcal N\!\left(\mu,\frac{\sigma^2}{n}\right).
\]
For Gaussian observations this distribution is exact for every \(n\),
by Chapter~\ref{chap:joint-distributions}. For exponential observations
it is an approximation to a Gamma density. Figure~\ref{fig:ch4-clt}
shows the rescaled densities: the narrowing has been removed, allowing
us to see the change in shape.

\begin{figure}[H]
  \centering
  \begin{tikzpicture}
\begin{groupplot}[
  group style={group size=2 by 2,horizontal sep=1.15cm,vertical sep=1.6cm},
  width=0.44\textwidth,height=4.9cm,
  axis lines=left,xmin=-4,xmax=5,ymin=0,ymax=1.12,
  xtick={-4,-2,0,2,4},ytick={0,0.5,1},
  xlabel={Standardized average \(z\)},ylabel={Density},
  tick label style={font=\small},label style={font=\small},
  title style={font=\small,text=CourseDark},
  grid=major,grid style={black!8},
]
\nextgroupplot[title={\(n=1\)}]
  \addplot[CourseBlue,very thick] table[x=z,y=exact] {figures/data/ch4-clt-n1.dat};
  \addplot[orange!85!black,thick,dashed] table[x=z,y=normal] {figures/data/ch4-clt-n1.dat};
  \addplot[CourseGray,thin,densely dotted] coordinates {(-1,0) (-1,1.04)};
  \node[font=\footnotesize,text=CourseGray,anchor=north east,fill=white,inner sep=2pt]
    at (axis cs:-1.15,0.86) {\(z\geq-1\)};
\nextgroupplot[title={\(n=4\)}]
  \addplot[CourseBlue,very thick] table[x=z,y=exact] {figures/data/ch4-clt-n4.dat};
  \addplot[orange!85!black,thick,dashed] table[x=z,y=normal] {figures/data/ch4-clt-n4.dat};
  \addplot[CourseGray,thin,densely dotted] coordinates {(-2,0) (-2,0.76)};
  \node[font=\footnotesize,text=CourseGray,anchor=south west,fill=white,inner sep=2pt]
    at (axis cs:-2,0.77) {\(z\geq-2\)};
\nextgroupplot[title={\(n=25\)}]
  \addplot[CourseBlue,very thick] table[x=z,y=exact] {figures/data/ch4-clt-n25.dat};
  \addplot[orange!85!black,thick,dashed] table[x=z,y=normal] {figures/data/ch4-clt-n25.dat};
  \node[font=\footnotesize,text=CourseGray,anchor=north,align=center]
    at (rel axis cs:0.5,0.92) {Support begins at \(-5\)\\(outside this window)};
\nextgroupplot[title={\(n=100\)}]
  \addplot[CourseBlue,very thick] table[x=z,y=exact] {figures/data/ch4-clt-n100.dat};
  \addplot[orange!85!black,thick,dashed] table[x=z,y=normal] {figures/data/ch4-clt-n100.dat};
  \node[font=\footnotesize,text=CourseGray,anchor=north,align=center]
    at (rel axis cs:0.5,0.92) {Support begins at \(-10\)\\(outside this window)};
\end{groupplot}
\draw[CourseBlue,very thick] ([xshift=-4.7cm,yshift=0.85cm]group c1r1.north east) -- ++(0.55cm,0)
  node[right,font=\small,text=CourseDark] {Exact density};
\draw[orange!85!black,thick,dashed] ([xshift=0.35cm,yshift=0.85cm]group c1r1.north east) -- ++(0.55cm,0)
  node[right,font=\small,text=CourseDark] {Standard normal};
\end{tikzpicture}

  \caption{Exact densities of
  \(Z_n=\sqrt n(\bar T_n-5)/5\), compared with the standard Gaussian.
  All panels use the same axes. The original waits remain exponential;
  it is the standardized average whose shape becomes nearly Gaussian.
  The lower support endpoint is \(-\sqrt n\), so the small-sample
  curves are visibly asymmetric.}
  \label{fig:ch4-clt}
\end{figure}

\subsection{A bound, an approximation, and an exact answer}

For one hundred waits, the one-minute error tolerance is two standard
errors. The Gaussian approximation gives
\begin{align*}
  \Prob(|\bar T_{100}-5|>1)
  &=\Prob(|Z_{100}|>2)\\
  &\approx 2[1-\Phi(2)]\approx0.04550.
\end{align*}
The exact Gamma distribution in \eqref{eq:exponential-sample-mean}
lets us check it. Let \(G_n\) be the CDF of
\(\operatorname{Gamma}(n,0.2n)\). The count--waiting-time identity
\eqref{eq:arrivals-counts} gives, for \(t\ge0\),
\[
  G_n(t)=1-e^{-0.2nt}\sum_{j=0}^{n-1}\frac{(0.2nt)^j}{j!}.
\]
Numerically,
\[
  G_{100}(4)\approx0.017108,
  \qquad 1-G_{100}(6)\approx0.027864.
\]
The two tails are unequal because the exact sampling distribution
remains somewhat skewed. Their sum is about \(0.04497\).

\begin{center}
\begin{tabularx}{\textwidth}{@{}lclY@{}}
\toprule
Method & Value & Meaning & Information used\\
\midrule
Chebyshev & \(0.25000\) & Upper bound & Independence and variance\\
CLT & \(0.04550\) & Approximation & Gaussian approximation to the average\\
Gamma & \(0.04497\) & Exact probability & Full exponential model\\
\bottomrule
\end{tabularx}
\end{center}
The approximation is close here, but it is not a guarantee. Chebyshev
is conservative here, but remains valid without the exponential shape
assumption. Figure~\ref{fig:ch4-tails} compares all three across sample
sizes rather than only at one hundred.

\begin{figure}[H]
  \centering
  \begin{tikzpicture}
\begin{axis}[
  width=0.95\textwidth,height=6.9cm,
  axis lines=left,xmin=1,xmax=200,ymode=log,ymin=0.003,ymax=1.2,
  xlabel={Number of observations \(n\)},
  ylabel={\(\Prob(|\overline T_n-5|>1)\) (log scale)},
  xtick={1,50,100,150,200},ytick={0.01,0.1,1},yticklabels={0.01,0.1,1},
  tick label style={font=\small},label style={font=\small},
  grid=major,grid style={black!8},
  legend style={draw=none,font=\small,at={(0.5,1.04)},anchor=south,
    legend columns=3,column sep=0.6em},
]
  \addplot[CourseGray,thin,densely dotted,forget plot] coordinates {(100,0.003) (100,1)};
  \addplot[CourseBlue,very thick] table[x=n,y=exact] {figures/data/ch4-tail-comparison.dat};
  \addlegendentry{Exact Gamma}
  \addplot[orange!85!black,very thick,dashed] table[x=n,y=chebyshev] {figures/data/ch4-tail-comparison.dat};
  \addlegendentry{Chebyshev bound}
  \addplot[teal!80!black,thick,densely dotted] table[x=n,y=clt] {figures/data/ch4-tail-comparison.dat};
  \addlegendentry{CLT approximation}
  \node[font=\small,text=CourseGray,anchor=north west,fill=white,inner sep=2pt]
    at (axis cs:102,0.9) {\(n=100\)};
\end{axis}
\end{tikzpicture}

  \caption{Probability that the average wait misses five minutes by
  more than one minute. The exact Gamma probability and CLT
  approximation are close at moderate sample sizes in this example;
  the Chebyshev curve is an upper bound, capped at one. The vertical
  axis is logarithmic so small probabilities remain distinguishable.}
  \label{fig:ch4-tails}
\end{figure}

There is no universal sample size, such as thirty, at which the
Gaussian approximation becomes reliable for every distribution and
every probability. Strong skewness, rare extreme values, and far-tail
questions can require much larger samples. The theorem states
convergence, not a numerical error guarantee at a chosen finite \(n\).
In particular, approximation of an extremely small probability may
have poor relative accuracy even when the CDF approximation is good
in absolute terms.

\section{The limits of averaging}
\label{sec:averaging-failures}

\subsection{A shared disturbance does not average away}

The sensor example in Section~\ref{sec:random-linear-transforms}
showed that correlated errors resist averaging. To see the consequence
as sample size grows, suppose all observations have variance
\(\sigma^2\) and each distinct pair has correlation
\(\rho\in[0,1]\). The covariance formula
\eqref{eq:variance-quadratic-form} gives
\begin{align}
  \Var(\bar X_n)
  &=\frac{n\sigma^2+n(n-1)\rho\sigma^2}{n^2}\notag\\
  &=\sigma^2\left(\rho+\frac{1-\rho}{n}\right).
  \label{eq:correlated-mean-variance}
\end{align}
There are \(n\) diagonal terms and \(n(n-1)\) off-diagonal terms.
If \(\rho>0\), the variance approaches \(\rho\sigma^2\), not zero.

A model that realizes this behavior is
\[
  X_i=\mu+B+\varepsilon_i,
\]
where \(B\) is a centered calibration error shared by the whole
dataset. A fresh \(B\) is drawn for a new independently calibrated
dataset, but it stays fixed within that dataset's \(n\) readings.
The \(\varepsilon_i\) are centered independent reading
errors, independent of \(B\). Set
\(\Var(B)=\rho\sigma^2\) and
\(\Var(\varepsilon_i)=(1-\rho)\sigma^2\). Then
\[
  \bar X_n=\mu+B+\frac1n\sum_{i=1}^n\varepsilon_i.
\]
The independent reading errors average toward zero, leaving the
shared error. Repeating readings with the same calibration cannot
resolve it. Independent calibrations would provide different information.

\begin{figure}[H]
  \centering
  \begin{tikzpicture}
\begin{axis}[
  width=0.95\textwidth,height=6.5cm,
  axis lines=left,xmode=log,log basis x=10,xmin=1,xmax=1000,ymin=0,ymax=1.05,
  xlabel={Number of observations \(n\) (log scale)},
  ylabel={Standard error / individual SD},
  xtick={1,10,100,1000},xticklabels={1,10,100,1000},
  ytick={0,0.2,0.4,0.6,0.8,1},
  tick label style={font=\small},label style={font=\small},
  grid=major,grid style={black!8},
  legend style={draw=none,font=\small,at={(0.5,1.04)},anchor=south,
    legend columns=2,column sep=1.1em},
]
  \addplot[CourseBlue,very thick] table[x=n,y=independent] {figures/data/ch4-dependent-averages.dat};
  \addlegendentry{Independent observations}
  \addplot[orange!85!black,very thick,dashed] table[x=n,y=dependent] {figures/data/ch4-dependent-averages.dat};
  \addlegendentry{Pairwise correlation \(\rho=0.1\)}
  \addplot[CourseGray,thin,densely dotted,forget plot] coordinates {(1,0.316227766) (1000,0.316227766)};
  \node[font=\small,text=CourseDark,anchor=south west,fill=white,inner sep=3pt]
    at (axis cs:24,0.39) {Nonzero limit: \(\sqrt{0.1}\approx0.316\)};
\end{axis}
\end{tikzpicture}

  \caption{Standard error relative to the single-observation standard
  deviation. Under independence it decreases as \(1/\sqrt n\).
  A shared component giving pairwise correlation \(0.1\) leaves a
  limiting value \(\sqrt{0.1}\). Both curves use the same marginal
  variance; only the dependence changes.}
  \label{fig:ch4-dependence}
\end{figure}

For \(n=100\) and \(\rho=0.1\), the variance is
\(0.109\sigma^2\), versus \(0.01\sigma^2\) under independence.
Equating this variance to \(\sigma^2/n_{\mathrm{eff}}\) gives
\[
  n_{\mathrm{eff}}=\frac{n}{1+(n-1)\rho}\approx9.17.
\]
In this specific equal-correlation model, one hundred readings have
the same mean variance as roughly nine independent readings. This
\emph{effective sample size} is a variance comparison, not a claim
that the entire datasets have identical distributions. Other forms
of dependence can still permit an LLN or a CLT, but require their
own assumptions and variance calculation.

\subsection{Finite mean and finite variance are different requirements}

The Student distributions in Section~\ref{sec:student} already supply
the examples we need. For i.i.d. observations with distribution
\(t_{3/2}\), the absolute first moment is finite but the variance is
infinite. The general i.i.d. law of large numbers therefore applies:
the sample mean converges in probability to zero. The finite-variance
Chebyshev argument and the classical CLT above do not apply. There is
no finite \(\sigma\) to insert into \(\sigma/\sqrt n\).

Standard Cauchy observations (\(t_1\)) have an exact averaging property:
\begin{equation}
  X_1,\ldots,X_n\stackrel{\mathrm{i.i.d.}}{\sim}t_1
  \quad\Longrightarrow\quad
  \bar X_n\sim t_1\quad\text{for every }n.
  \label{eq:cauchy-average}
\end{equation}
Here a sample average is a finite arithmetic calculation for any
realized finite dataset, even though its population expectation does
not exist. Averaging independent Cauchy observations leaves their
distribution unchanged. In particular,
\[
  \Prob(|\bar X_n|>1)=\frac12
\]
for every \(n\), rather than tending to zero. Zero is the distribution's
center of symmetry, not its mean. Figure~\ref{fig:ch4-heavy-tails}
contrasts this with Gaussian averaging.

\begin{figure}[H]
  \centering
  \begin{tikzpicture}
\begin{groupplot}[
  group style={group size=2 by 1,horizontal sep=1.25cm},
  width=0.44\textwidth,height=5.8cm,
  axis lines=left,xmin=-4,xmax=4,ymin=0,
  xlabel={Sample average},ylabel={Density},xtick={-4,-2,0,2,4},
  tick label style={font=\small},label style={font=\small},
  title style={font=\small,text=CourseDark},
  grid=major,grid style={black!8},
]
\nextgroupplot[ymax=2.15,ytick={0,0.5,1,1.5,2},title={Standard Gaussian}]
  \addplot[CourseBlue,very thick] table[x=x,y=n1] {figures/data/ch4-heavy-tail-means.dat};
  \addplot[orange!85!black,very thick,dashed] table[x=x,y=n4] {figures/data/ch4-heavy-tail-means.dat};
  \addplot[teal!80!black,very thick,densely dotted] table[x=x,y=n25] {figures/data/ch4-heavy-tail-means.dat};
  \node[font=\footnotesize,text=CourseDark,anchor=north west,align=left]
    at (rel axis cs:0.03,0.96) {\(n=25\):\\narrowest};
\nextgroupplot[ymax=0.36,ytick={0,0.1,0.2,0.3},title={Standard Cauchy}]
  \addplot[CourseDark,very thick] table[x=x,y=cauchy] {figures/data/ch4-heavy-tail-means.dat};
  \node[font=\footnotesize,text=CourseDark,anchor=north west,align=left,
    fill=white,inner sep=2pt]
    at (rel axis cs:0.01,0.96) {Same density\\for every \(n\)};
\end{groupplot}
\draw[CourseBlue,very thick] ([xshift=-3.7cm,yshift=0.8cm]group c1r1.north east) -- ++(0.5cm,0)
  node[right,font=\small,text=CourseDark] {\(n=1\)};
\draw[orange!85!black,very thick,dashed] ([xshift=-1.6cm,yshift=0.8cm]group c1r1.north east) -- ++(0.5cm,0)
  node[right,font=\small,text=CourseDark] {\(n=4\)};
\draw[teal!80!black,very thick,densely dotted] ([xshift=0.5cm,yshift=0.8cm]group c1r1.north east) -- ++(0.5cm,0)
  node[right,font=\small,text=CourseDark] {\(n=25\)};
\end{tikzpicture}

  \caption{Exact densities of sample means from two centered symmetric
  populations. Gaussian means concentrate as \(n\) grows. For standard
  Cauchy observations the density is unchanged for every sample size,
  so the three Cauchy curves coincide. The panels have different
  vertical scales. Symmetry alone does not supply
  the moment assumptions needed for averaging to become reliable.}
  \label{fig:ch4-heavy-tails}
\end{figure}

The conclusions are distinct: infinite variance need not destroy the
LLN, and independence alone does not guarantee that a sample mean
concentrates. Check the assumptions for the particular theorem being
used, not just whether the data look roughly symmetric.

\section{Monte Carlo and averages in machine learning}
\label{sec:monte-carlo}

\subsection{Replace an integral by independent evaluations}

Suppose the quantity we need is \(I=\Expect[g(X)]\). When evaluating
the corresponding integral is difficult but generating observations
is feasible, draw \(X_1,\ldots,X_n\) independently from the target
distribution and compute
\begin{equation}
  \widehat I_n=\frac1n\sum_{i=1}^n g(X_i).
  \label{eq:monte-carlo}
\end{equation}
This is \emph{Monte Carlo estimation}. The earlier theorems now apply
to \(Y_i=g(X_i)\). If \(\Expect[|g(X)|]<\infty\), the LLN gives
convergence to \(I\). If also
\(0<\tau^2=\Var(g(X))<\infty\), then
\[
  \operatorname{SE}(\widehat I_n)=\frac{\tau}{\sqrt n},\qquad
  \frac{\sqrt n(\widehat I_n-I)}{\tau}
  \xrightarrow{d}\mathcal N(0,1).
\]
The relevant variance is that of the \emph{evaluated quantity}
\(g(X)\), not necessarily that of \(X\). Even a heavy-tailed input
can give a bounded output, such as an event indicator, whose variance
is finite.

Use the response score from Section~\ref{sec:expectation} as a case
with a known answer. For \(T\sim\operatorname{Exp}(0.2)\) and
\(g(t)=e^{-t}\), with time expressed in minutes, the earlier moment
calculation gives
\[
  I=\frac16,\qquad
  \Expect[g(T)^2]=\frac1{11},\qquad
  \tau^2=\frac1{11}-\frac1{36}=\frac{25}{396}.
\]
The coefficient of \(t\) in the exponent is one per minute, making
the exponent dimensionless. We can compare simulated estimates
against \(1/6\) to check the numerical method before using it for
an expectation with no closed-form answer.

\subsection{Precision has a computational cost}

Because \(\Expect[\widehat I_n]=I\), its root mean squared error over
independent datasets is
\begin{equation}
  \operatorname{RMSE}(\widehat I_n)
  =\sqrt{\Expect[(\widehat I_n-I)^2]}
  =\frac{\tau}{\sqrt n}.
  \label{eq:mc-rmse}
\end{equation}
For the score, \(\tau\approx0.25126\); thus \(n=1000\) gives
RMSE about \(0.00795\). Multiplying the number of evaluations by
four halves this error scale. Reducing it by a factor of ten requires
one hundred times as many evaluations, if the variance stays the same.

For a target RMSE of \(0.005\), a sufficient integer sample size in
this example is
\[
  n\ge\left\lceil\frac{25/396}{0.005^2}\right\rceil=2526,
\]
where \(\lceil\cdot\rceil\) rounds upward. An RMSE target is not a
bound on the error in every run, nor does it specify a failure
probability. The two are different accuracy requirements.

\begin{figure}[H]
  \centering
  \begin{tikzpicture}
\begin{axis}[
  width=0.95\textwidth,height=6.9cm,
  axis lines=left,xmode=log,ymode=log,log basis x=10,log basis y=10,
  xmin=10,xmax=10000,ymin=0.002,ymax=0.12,
  xlabel={Number of independent draws \(n\) (log scale)},
  ylabel={Root mean square error (log scale)},
  xtick={10,100,1000,10000},xticklabels={10,100,1000,10000},
  ytick={0.003,0.01,0.03,0.1},yticklabels={0.003,0.01,0.03,0.1},
  tick label style={font=\small},label style={font=\small},
  grid=major,grid style={black!8},
  legend style={draw=none,font=\small,at={(0.5,1.04)},anchor=south,
    legend columns=2,column sep=1.1em},
]
  \addplot[CourseBlue,very thick] table[x=n,y=theory] {figures/data/ch4-monte-carlo-error.dat};
  \addlegendentry{Theory: \(\sqrt{25/(396n)}\)}
  \addplot[orange!85!black,only marks,mark=*,mark size=2.1pt,
    mark options={fill=white,line width=0.9pt}]
    table[x=n,y=simulated] {figures/data/ch4-monte-carlo-error.dat};
  \addlegendentry{Simulation (2,000 repetitions)}
  \draw[CourseGray,thick,dashed] (axis cs:100,0.07) -- (axis cs:400,0.07)
    -- (axis cs:400,0.035);
  \node[font=\small,text=CourseGray,anchor=south] at (axis cs:200,0.073) {\(4\times n\)};
  \node[font=\small,text=CourseGray,anchor=west,fill=white,inner sep=2pt]
    at (axis cs:420,0.0495) {\(\tfrac12\times\) error};
\end{axis}
\end{tikzpicture}

  \caption{Monte Carlo error for the response score \(e^{-T}\).
  The theoretical curve is the exact RMSE \(\sqrt{25/396}/\sqrt n\).
  The simulated points estimate RMSE from repeated independent runs;
  they are not the absolute error of one running estimate. Both axes
  are logarithmic, so the \(n^{-1/2}\) rate is a straight line.}
  \label{fig:ch4-monte-carlo}
\end{figure}

When \(\tau\) is unknown, compute the sample standard deviation
\(S_Y\) of the values \(g(X_i)\) and report \(S_Y/\sqrt n\) as an
estimated standard error. Increasing \(n\) does not fix a wrong
sampling distribution or a faulty implementation of \(g\).

\subsection{A test-set average estimates a population quantity}

Fix a prediction rule \(f\) before examining the evaluation data.
Let \((X_i,Y_i)\) be i.i.d. test observations, independent of any
data used to choose \(f\). For a per-observation loss
\(L_i=\ell(f(X_i),Y_i)\), the population quantity and its empirical
average are
\[
  R(f)=\Expect[\ell(f(X),Y)],\qquad
  \widehat R_n(f)=\frac1n\sum_{i=1}^n\ell(f(X_i),Y_i).
\]
This is the same Monte Carlo construction. Conditional on the fixed
model, the variation comes from which evaluation examples happen to
be sampled. With finite loss variance,
\[
  \operatorname{SE}(\widehat R_n(f))
  =\frac{\sqrt{\Var(\ell(f(X),Y))}}{\sqrt n}.
\]
The squared-error interpretation from
Section~\ref{sec:conditional-moments} supplies one choice of loss;
classification error indicators supply another. The new point is
that a reported average itself fluctuates across test sets.

For a fixed classifier with population accuracy \(0.8\), four hundred
independent test examples give standard error
\[
  \sqrt{\frac{0.8(0.2)}{400}}=0.02.
\]
This describes test-set sampling variation,
not variation caused by retraining the model. Selecting the model
with the best score on that same test set breaks the fixed-model
premise. More observations from an unrepresentative test population
also do not establish performance on the intended population.

\subsection{What increasing the batch size changes}

At a fixed model, average the losses of \(b\) independent examples.
If each has variance \(\tau_L^2\), the batch-average variance is
\(\tau_L^2/b\). Thus a fourfold batch-size increase halves the
standard deviation of that average. This is an averaging statement,
not a claim that training becomes four times better or faster.

For a fixed finite dataset, independent sampling of indices
\emph{with replacement} gives the same rule conditional on that
dataset, using its empirical loss variance. Sampling without
replacement gives a different variance, and correlated examples need
the covariance terms from \eqref{eq:correlated-mean-variance} or the
general quadratic-form rule. The number of rows in a batch is not
enough to determine its information content.

The three tools answer different questions. Chebyshev bounds an error
probability at a finite sample size. The LLN establishes convergence
under growing sample size. The CLT describes rescaled fluctuations
and supports approximations. A reliable use of averages keeps those
claims, and their assumptions, separate.
\endgroup
